% =============================================================================
% chapter1.tex — Глава 1. Анализ предметной области
% =============================================================================

\section{АНАЛИЗ ПРЕДМЕТНОЙ ОБЛАСТИ И СУЩЕСТВУЮЩИХ РЕШЕНИЙ}

\subsection{Проблема поддержания актуальной документации в корпоративных проектах}

\todo{Описать проблему деградации документации: статистика (Forward, Lethbridge: 60--70\% документации устаревает за несколько месяцев~\cite{forward2002relevance}), причины (хотфиксы, устные договорённости, давление бизнеса, темпы CI/CD), последствия (технический долг, сложность онбординга, риски при модернизации). Привести реальный пример: в проекте doc-generator за один спринт было внесено 38 коммитов~--- ни один не затронул документацию.}

\subsection{Обзор существующих подходов к автоматической генерации документации}

\subsubsection{Code summarization~--- генерация кратких описаний из кода}

\todo{CodeBERT~\cite{feng2020codebert}, CodeT5~\cite{wang2021codet5}, GraphCodeBERT~\cite{guo2021graphcodebert}: возможности (генерация docstring, перевод код→текст) и ограничения (контекст одной функции, невозможность учёта вызовных цепочек, отсутствие знаний о межмодульных зависимостях). Локальный контекст: модели работают на уровне одного метода/класса, не видят интеграционных связей (HTTP, Kafka).}

\subsubsection{LLM в IDE (IDE-plugins и AI copilots)}

\todo{GitHub Copilot~\cite{nguyen2022copilot}, JetBrains AI Assistant, TabNine: сильные стороны (контекст открытого файла, автодополнение, генерация docstring на лету) и ограничения (ограниченное контекстное окно~--- 8K-128K токенов, зависимость от облака, проблемы конфиденциальности корпоративного кода~\cite{vaithilingam2022usability}, невозможность целостного описания сервиса, отсутствие архитектурного контекста).}

\subsubsection{Автоматическая API-документация (Swagger/OpenAPI + AI)}

\todo{Генерация из аннотаций (@GetMapping, @PostMapping и др.): Swagger UI, springdoc-openapi~\cite{westhofen2022openapi,lyu2024openapi}. LLM для автозаполнения описаний. Ограниченность: покрывает только интерфейсный REST-слой, не описывает бизнес-логику, вызовные цепочки, Kafka-интеграции, Camel-маршруты, Circuit Breaker паттерны.}

\subsubsection{RAG над кодовой базой}

\todo{Retrieval-Augmented Generation~\cite{siddiq2024rag}: схема работы (embedding→vector store→retrieval→LLM), SourceGraph Cody, Repo-RAG~\cite{shrivastava2024reporag}. Качество поиска: проблемы (Retrieval Quality: наивный чанкинг файлов теряет структуру, <<потеря в середине>>~\cite{liu2023lost}). Без графовой модели: RAG не видит вызовные цепочки, не знает о CALLS\_CODE/CALLS\_HTTP/PRODUCES/CONSUMES связях, ответ фрагментарен.}

\subsubsection{Code Graph + LLM~--- графовая репрезентация системы}

\todo{Граф знаний кода~\cite{xu2024codegraph,peng2024structural}: узлы (классы, методы, эндпоинты, топики), рёбра (CALLS\_CODE, CALLS\_HTTP, PRODUCES, CONSUMES), AST/PSI. Возможность прослеживания путей выполнения: METHOD→CALLS\_CODE→METHOD→CALLS\_HTTP→ENDPOINT. Сложность построения (необходимость полного classpath для разрешения типов), но уникальное преимущество~--- системный контекст: один метод описывается вместе со всеми его зависимостями.}

\subsection{Сравнительный анализ подходов}

\todo{Таблица сравнения 5 подходов по 7 критериям:

\begin{center}
\begin{tabular}{|l|c|c|c|c|c|}
\hline
Критерий & Code Sum & IDE Copilot & Swagger & RAG & Graph+LLM \\
\hline
Контекст метода & + & + & -- & +/-- & + \\
Контекст сервиса & -- & -- & +/-- & +/-- & + \\
Контекст системы & -- & -- & -- & -- & + \\
Интеграции & -- & -- & REST only & -- & HTTP/Kafka/Camel \\
Конфиденциальность & +/-- & -- & + & +/-- & + (локально) \\
Сложность внедрения & низкая & низкая & низкая & средняя & высокая \\
Автоактуальность & -- & -- & + & -- & + \\
\hline
\end{tabular}
\end{center}

Обосновать оценки со ссылками на источники. Вывод: ни один из существующих подходов не обеспечивает одновременно контекст системы, конфиденциальность и автоактуальность. Предлагаемый подход Graph+RAG+LLM объединяет преимущества.}

\subsection{Постановка задачи и обоснование выбранного подхода}

\todo{Обоснование интеграции графовой модели + RAG + LLM. Формулировка контура обработки:

GIT (clone/pull) → PSI-парсинг Kotlin → Code Graph (24 типа узлов, 24 типа связей) → Семантическое чанкование (PerNodeChunkStrategy) → Embeddings (bge-m3, 1024-dim) → pgvector → Двухступенчатая генерация (Coder→Talker→Digest) → Документация

Параллельно: ASM-анализ библиотек → обнаружение интеграционных точек (HTTP, Kafka, Camel)

Интерактивный режим: RAG с графовым расширением контекста (8 шагов обработки запроса)

Визуализация: Cytoscape.js (Single-App + Cross-App режимы)

Открытые вопросы: выбор локальных LLM для конфиденциальности (Ollama + Qwen 2.5), стратегия чанкования (один узел — один чанк), обеспечение отказоустойчивости (Circuit Breaker, Retry, Bulkhead).}
